\documentclass[aspectratio=169,xcolor={dvipsnames,table}]{beamer}
\usepackage[no-math,deluxe,expert,haranoaji]{luatexja-preset}
\usepackage{luatexja-otf}
\renewcommand{\kanjifamilydefault}{\gtdefault}
\renewcommand{\emph}[1]{{\upshape\bfseries #1}}
\usetheme{metropolis}
\metroset{block=fill}
\setbeamertemplate{navigation symbols}{}
\setbeamertemplate{blocks}[rounded][shadow=false]
\usecolortheme[rgb={0.7,0.2,0.2}]{structure}
%%%%%%%%%%%%%%%%%%%%%%%%%%%
\usepackage{media9}
\usepackage[absolute,overlay]{textpos}
%\usepackage[grid=true,gridcolor=Maroon,subgridcolor=gray,gridunit=pt,texcoord]{eso-pic} %場所決めのためのgrid表示
%%%%%%%%%%%%%%%%%%%%%%%%%%%
%% さまざまなアイコン
%%%%%%%%%%%%%%%%%%%%%%%%%%%
\usepackage{fontawesome}
%\usepackage{figchild}
\usepackage{twemojis}
\usepackage{utfsym}
\usepackage{bclogo}
\usepackage{marvosym}
\usepackage{fontmfizz}
\usepackage{pifont}
\usepackage{phaistos}
\usepackage{worldflags}
%%%%%%%%%%%%%%%%%%%%%%%%%%%
\usepackage{tikz}
\usetikzlibrary{backgrounds}
\usetikzlibrary{tikzmark}
\usepackage{tcolorbox}
\usepackage{tikzpeople}
\usepackage{tikzducks}
\usepackage{tikzlings}
\usepackage{circledsteps}
\usepackage{xcolor}
\usepackage{amsmath}
\usepackage{booktabs}
\usepackage{chronology}
\usepackage{signchart}
\usepackage{tipa}
\usepackage{pxrubrica}
%%%%%%%%%%%%%%%%%%%%%%%%%%%
%% 場合分け
\usepackage{cases}
%%%%%%%%%%%%%%%%%%%%%%%%%%%
% \myAnch{<名前>}{<色>}{<テキスト>}
% 指定のテキストを指定の色の四角枠で囲み, 指定の名前をもつTikZの
% ノードとして出力する. 図には remeber picture 属性を付けている
% ので外部から参照可能である.
\newcommand*{\myAnch}[3]{%
  \tikz[remember picture,baseline=(#1.base)]
    \node[draw,rectangle,#2] (#1) {\normalcolor #3};
}
%%%%%%%%%%%%%%%%%%%%%%%%%%%%
%% 音声リンク表示
\newcommand{\myaudio}[1]{\href{#1}{\faVolumeUp}}
%%%%%%%%%%%%%%%%%%%%%%%%%%%
% \myEmph コマンドの定義
%\newcommand{\myEmph}[3]{%
%    \textbf<#1>{\color<#1>{#2}{#3}}%
%}
\usepackage{xparse} % xparseパッケージの読み込み
\NewDocumentCommand{\myEmph}{O{} m m}{%
    \def\argOne{#1}%
    \ifx\argOne\empty
        \textbf{\color{#2}{#3}}% オプション引数が省略された場合
    \else
        \textbf<#1>{\color<#1>{#2}{#3}}% オプション引数が指定された場合
    \fi
}
%%%%%%%%%%%%%%%%%%%%%%%%%%%
%% 文末の上昇イントネーション記号 \myRisingPitch
%% 通常のイントネーション \myDownwardPitch
%% https://note.com/dan_oyama/n/n8be58e8797b2
%%%%%%%%%%%%%%%%%%%%%%%%%%%
\newcommand{\myRisingPitch}{
\begin{tikzpicture}[scale=0.3,baseline=0.3]
\draw[->,>=stealth] (0,0) to[bend right=45] (1,1);
\end{tikzpicture}
}
\newcommand{\myDownwardPitch}{
\begin{tikzpicture}[scale=0.3,baseline=0.3]
\draw[->,>=stealth] (0,1) to[bend left=45] (1,0);
\end{tikzpicture}
}
%%%%%%%%%%%%%%%%%%%%%%%%%%%
\title{English is fun.}
\subtitle{I have known him for eight years.}
  \author{}
\institute[]{}
\date[]

%%%%%%%%%%%%%%%%%%%%%%%%%%%%
%% TEXT
%%%%%%%%%%%%%%%%%%%%%%%%%%%%
\begin{document}
\begin{frame}[plain]
  \titlepage
\end{frame}

\section*{授業の流れ}
\begin{frame}[plain]
  \frametitle{授業の流れ}
  \tableofcontents
\end{frame}

%%%%%%%%%%%%%%%%%%%%%%%%
\section*{現在完了とは（復習）}
%%%%%%%%%%%%%%%%%%%%%%%%%%%%
\begin{frame}[plain]{現在完了とは(復習)}
 \begin{enumerate}
 \item Jane \textcolor{NavyBlue}{\bfseries has stayed} in London for six days.\visible<6->{（継続）}
 \item I \textcolor{NavyBlue}{\bfseries have watched} the movie three times.\visible<7->{（経験）}
 \item Bob \textcolor{NavyBlue}{\bfseries has lost} his bag.\visible<8->{（結果）}
\end{enumerate}

\vspace{30pt}

 \begin{block}<2->{現在完了の基本}
\small
\begin{itemize}\setbeamertemplate{items}[square]
 \item<3->  現在完了形\,\Circled[fill color=white]{\,\textbf{have} $+$ 過去分詞\,}\,は「過去と現在にまたがる表現」です\\
\hfill\visible<4->{{\scriptsize 主語が三人称単数のときは$\textcolor{NavyBlue}{\text{\bfseries has} + \text{過去分詞}}$}}
 \item \visible<5->{過去分詞は\kenten{受け身}とならんで\kenten{完了}の意味をあらわします}%
\end{itemize}
      \end{block}
\hfill{\tiny 0144}\,{\scriptsize \myaudio{./audio/012_have_pp_keizoku_01.mp3}}

\begin{textblock*}{0.4\linewidth}(365pt,50pt)
\visible<1->{\begin{tikzpicture}
\duck[signpost=\scalebox{0.3}{
\parbox{2.5cm}{\color{black}\centering
{\Large 現在完了$=$\\過去$+$現在}}},
signcolour=brown!70!gray,
signback=white!80!brown,
graduate=gray!20!black,
tassel=red!70!black,
laughing,
%speech={\tiny 反復練習！}
]
\end{tikzpicture}}
\end{textblock*}
\end{frame}
%%%%%%%%%%%%%%%%%%%%%%%%
\section{継続をあらわす現在完了}
%%%%%%%%%%%%%%%%%%%%%%
\begin{frame}[plain,t]{現在完了--継続--}
 
\visible<1->{\begin{enumerate}
 \item $\left\{\begin{tabular}{rl}
(A)& Jane \textcolor{Maroon}{\bfseries stayed} in London two days ago.\hspace{80pt}{\scriptsize stay \textipa{/st\'eI/} 滞在する}\\
(B)& Jane \textcolor{NavyBlue}{\bfseries has stayed} in London for six days.\hspace{69pt}{\scriptsize for \textipa{/f\textrhookschwa /} ～の間}\\
\end{tabular}
\right.$
\end{enumerate}}

\visible<2->{\signchart[width=10,height=.5]{,,,,,{\textcolor{Maroon}{\textbf{stayed}}},,今}{,,,}}
\visible<3->{\signchart[width=10,height=.5]{,6日前,,,{\textcolor{NavyBlue}{\textbf{has stayed}}},,,今}{,,,}}

\begin{tikzpicture}[overlay]
 %\draw[gray!50] (0,0) grid (12,5);
 %\fill[ForestGreen!70,opacity=.5] (11.57,3.75) circle [radius=.25];
\visible<2->{ \fill[Maroon!70,opacity=.5] (8.675,2.66) circle [radius=.25];}
\visible<3->{ \draw[NavyBlue!70,line width=6pt,opacity=.7] (4.2,1.21) -- (10.9,1.21);}
\visible<3->{\node (kikan) at (7.5,0.5) {6日間};}
\visible<3->{\draw (4.2,1.2) to[out=-30,in=180] (kikan.west);}
\visible<3->{\draw (kikan.east) to[out=0,in=210] (10.9,1.2);}
\end{tikzpicture}

\visible<4->{{\scriptsize 主語が三人称単数のときは has $+$ 過去分詞}}

\vspace{-5pt}

\visible<5->{{\scriptsize \textbf{for} six daysは\kenten{期間}を表します}}%
\hfill{\tiny 0120}\,{\scriptsize \myaudio{./audio/012_have_pp_keizoku_02.mp3}}

\begin{textblock*}{0.4\linewidth}(385pt,90pt)
\visible<2->{\begin{tikzpicture}
\duck[signpost=\scalebox{0.3}{
\parbox{2.5cm}{\color{black}\centering
{\Large 1目盛り\\1日}！}},
signcolour=brown!70!gray,
signback=white!80!brown,
graduate=gray!20!black,
tassel=red!70!black,
laughing,
%speech={\tiny 反復練習！}
]
\end{tikzpicture}}
\end{textblock*}
\end{frame}
%%%%%%%%%%%%%%%%%%%%%%
\begin{frame}[plain,t]{継続}
\visible<1->{\begin{enumerate}\setcounter{enumi}{1}
 \item $\left\{\begin{tabular}{rl}
(A)& I \textcolor{Maroon}{\bfseries am} sick today.\hspace{160pt}{\scriptsize sick \textipa{/s\'Ik/} 病気だ、具合が悪い}\\
%(B)& I \textcolor{NavyBlue}{\bfseries have been} sick since yeaterday.\\
		(B)& I \textcolor{NavyBlue}{\bfseries have been} sick for six days.\hspace{95pt}{\scriptsize been \textipa{/bIn/} beの過去分詞}
\end{tabular}
\right.$
\end{enumerate}}


\visible<2->{\signchart[width=10,height=.5]{,,,,,,,{\textcolor{Maroon}{\textbf{am}}}}{,,,}}
%\visible<3->{\signchart[width=10,height=.5]{,,,,,,{\textcolor{NavyBlue}{have been}},}{,,,}}

\visible<3->{\signchart[width=10,height=.5]{,6日前,,,{\textcolor{NavyBlue}{\textbf{have been}}},,,今}{,,,}}

\begin{tikzpicture}[overlay]
 %\draw[gray!50] (0,0) grid (12,5);
 \visible<2-> {\fill[Maroon!70,opacity=.5] (10.9,2.7) circle [radius=.25];}
%\visible<3->{ \draw[NavyBlue!70,line width=7pt,opacity=.7] (9.8,2.65) -- (10.9,2.65);}
\visible<3->{ \draw[NavyBlue!70,line width=7pt,opacity=.7] (4.2,1.21) -- (10.9,1.21);}
\visible<3->{\node (kikan) at (7.5,0.5) {6日間};}
\visible<3->{\draw (4.2,1.2) to[out=-30,in=180] (kikan.west);}
\visible<3->{\draw (kikan.east) to[out=0,in=210] (10.9,1.2);}
\end{tikzpicture}

%\vspace{-20pt}

\visible<4->{{\scriptsize \textbf{for} six daysは\kenten{期間}を表します}}\hfill{\tiny 0112}\,{\scriptsize \myaudio{./audio/012_have_pp_keizoku_03.mp3}}

%\visible<5->{since 〜は「起点」を表す表現です。「〜から」}
\end{frame}
%%%%%%%%%%%%%%%%%%%%%%%%%%%%%%%
\begin{frame}[plain,t]{継続}
\visible<1->{\begin{enumerate}\setcounter{enumi}{2}
 \item $\left\{\begin{tabular}{rl}
(A)& I \textcolor{Maroon}{\bfseries know} him.\\
(B)& I \textcolor{NavyBlue}{\bfseries have known} him since 1996.%
\hspace{80pt}{\scriptsize since \textipa{/s\'Ins/} ～から、以来}
%(B)& I \textcolor{NavyBlue}{\bfseries have known} sick for six days.\hspace{3\zw}{\small known \textipa{/n\'oUn/} knowの過去分詞}
\end{tabular}
\right.$
\end{enumerate}}


\visible<2->{\signchart[width=10,height=.5]{,,,,,,,{\textcolor{Maroon}{\textbf{know}}}}{,,,}}
%\visible<3->{\signchart[width=10,height=.5]{,,,,,,{\textcolor{NavyBlue}{have been}},}{,,,}}

\visible<3->{\signchart[width=10,height=.5]{,1996,,,{\textcolor{NavyBlue}{\textbf{have known}}},,,2026}{,,,}}

\begin{tikzpicture}[overlay]
 %\draw[gray!50] (0,0) grid (12,5);
 \visible<2-> {\fill[Maroon!70,opacity=.5] (10.9,2.5) circle [radius=.25];}
%\visible<3->{ \draw[NavyBlue!70,line width=7pt,opacity=.7] (9.8,2.65) -- (10.9,2.65);}
\visible<3->{ \draw[NavyBlue!70,line width=7pt,opacity=.7] (4.2,1.21) -- (10.9,1.21);}
\visible<3->{\node (kikan) at (7.5,0.5) {30年間};}
\visible<3->{\draw (4.2,1.2) to[out=-30,in=180] (kikan.west);}
\visible<3->{\draw (kikan.east) to[out=0,in=210] (10.9,1.2);}
\end{tikzpicture}

%\vspace{-20pt}

%\visible<4->{for a weekは「期間」を表します。}
\visible<3->{\scriptsize \textbf{since} 〜は\kenten{起点}を表します。「〜から、以来」}\\%
\visible<3->{\scriptsize 1996の読み方：nineteen ninety-six}

\hfill{\tiny 0113}\,{\scriptsize \myaudio{./audio/012_have_pp_keizoku_04.mp3}}

\begin{textblock*}{0.4\linewidth}(385pt,90pt)
\visible<2->{\begin{tikzpicture}
\duck[signpost=\scalebox{0.3}{
\parbox{2.5cm}{\color{black}\centering
{\Large 1目盛り\\5年}！}},
signcolour=brown!70!gray,
signback=white!80!brown,
graduate=gray!20!black,
tassel=red!70!black,
laughing,
%speech={\tiny 反復練習！}
]
\end{tikzpicture}}
\end{textblock*}
\end{frame}
%%%%%%%%%%%%%%%%%%%%%%%%%%%%%%%%%%%%%%%
\section{forとsince}
%%%%%%%%%%%%%%%%%%%%%%%%%%%%%%%%%%%%%%
\begin{frame}[plain,t]{for \textipa{/f\textrhookschwa /}とsince \textipa{/s\'Ins/}}
 
\begin{enumerate}
 \item My parents {\bfseries have lived} in the same town \textcolor{NavyBlue}{\bfseries for} 40 years.%
       \hfill\visible<4->{{\scriptsize for \Circled{\,名詞\,}}}
 \item Billy {\bfseries has had} this car \textcolor{NavyBlue}{\bfseries for} five years.%
       \hfill\visible<4->{{\scriptsize for \Circled{\,名詞\,}}}
 \item I {\bfseries have been} sick \textcolor{Maroon}{\bfseries since} this morning.%
       \hfill\visible<6->{{\scriptsize since \Circled{\,名詞\,}}}
 \item I {\bfseries have known} him \textcolor{Maroon}{\bfseries since} he was a child.%
       \hfill\visible<7->{{\scriptsize since \Circled{\,S $+$ V\,}}}
\end{enumerate}

\vspace{30pt}

\begin{block}<2->{Topics for Today}\small
 \begin{itemize}\setbeamertemplate{items}[square]
  \item<3-> \textbf{for}は\kenten{期間}を、\textbf{since}は\kenten{起点}をあらわす
  \item<5-> \textbf{for}\tikzmark{for}\,\,\,\Circled[fill color=white]{\,名詞\,}\hspace{150pt}\visible<8->{{\scriptsize \tikzmark{prep}前置詞}}
  \item<8-> $\left\{
	\begin{tabular}{ll}
	\textbf{since}\tikzmark{since1}\,\,\,\Circled[fill color=white]{\,名詞\,}&\\
	\textbf{since}\tikzmark{since2}\,\,\,\Circled[fill color=white]{\,\textbf{S} $+$ \textbf{V}\,}&\hspace{106pt}\visible<10->{{\scriptsize \tikzmark{conj}接続詞}}
	\end{tabular}
	\right.$
 \end{itemize}
\end{block}
\vspace{-5pt}
\hfill{\tiny 0215}\,{\scriptsize \myaudio{./audio/012_have_pp_keizoku_05.mp3}}

\begin{tikzpicture}[remember picture,overlay]
 \visible<8->{\draw[<-,opacity=.6,line width=.8pt] ([xshift=-4pt,yshift=7pt]pic cs:for) to[out=22, in=180] ([xshift=-2pt, yshift=2pt] pic cs:prep);}
 \visible<9->{\draw[<-,opacity=.6,line width=.8pt] ([xshift=-2pt,yshift=8pt]pic cs:since1) to[out=45, in=180] ([xshift=-2pt, yshift=2pt] pic cs:prep);}
 \visible<10->{\draw[<-,opacity=.6,line width=.8pt] ([xshift=-2pt,yshift=-2pt]pic cs:since2) to[out=-15, in=180] ([xshift=-2pt, yshift=2pt] pic cs:conj);}
\end{tikzpicture}

\begin{textblock*}{0.4\linewidth}(310pt,133pt)
\visible<11->{\begin{tikzpicture}
\marmot[
speech={\scriptsize for{\tiny と}since},
bubblecolor=white,
signpost=\scalebox{0.5}{
\parbox{2.2cm}{\color{black}
\centering 意味の違いと\\品詞にも\\注意！}},
signcolour= brown!70!gray,
signback=white!80!brown
]
\end{tikzpicture}}
\end{textblock*}

\end{frame}
%%%%%%%%%%%%%%%%%%%%%%%%%%%%%%%%%%%%%%
\begin{frame}<1-11>[plain]{Exercises}

{\small あたえられた日本語の意味になるよう、(~~~~~~~~)内に適当な語を補いましょ}う%
\hfill{\tiny 0216}\,{\scriptsize \myaudio{./audio/012_have_pp_keizoku_05.mp3}}

\begin{enumerate}
 \item {\small 彼らは20年の間その家を所有しています。}\hfill{}{\scriptsize own \textipa{/\'oUn/} 所有する}\\
They \alt<1>{(~~\phantom{have}~~)}{(~~\textcolor{Maroon}{\bfseries have}~~)} owned this house for twenty years.%
       \hfill\visible<2->{{\scriptsize for \Circled{\,名詞\,}}}
 \item {\small ビリーは長い間、ジェーンを愛している。}%
       \hfill{}{\scriptsize practice \textipa{/pr\'\ae ktIs/} 練習する}\\
Billy \alt<1-2>{(~~\phantom{has}~~)}{(~~\textcolor{Maroon}{\bfseries has}~~)} loved Jane \alt<1-3>{(~~\phantom{for}~~)}{(~~\textcolor{Maroon}{\bfseries for}~~)} a long time.
 \item {\small わたしは今朝からずっと具合が悪い。}%
       \hfill{}{\scriptsize be sick: 病気だ、具合が悪い}\\
I have \alt<1-4>{(~~\phantom{been}~~)}{(~~\textcolor{Maroon}{\bfseries been}~~)} sick \alt<1-5>{(~~\phantom{since}~~)}{(~~\textcolor{Maroon}{\bfseries since}~~)} this morning.%
       \hfill\visible<6->{{\scriptsize since \Circled{\,名詞\,}}}
 \item {\small 彼が子供のころから、わたしは彼を知っています。}\\
I have \alt<1-6>{(~~\phantom{known}~~)}{(~~\textcolor{Maroon}{\bfseries known}~~)} him \alt<1-7>{(~~\phantom{since}~~)}{(~~\textcolor{Maroon}{\bfseries since}~~)} he was a child.%
       \hfill\visible<8->{{\scriptsize since \Circled{\,\,S $+$ V\,\,}}}
\end{enumerate}

\begin{block}<9->{Topics for Today}\small
\begin{itemize}\setbeamertemplate{items}[square]
 \item \visible<10->{{\bfseries for} 〜は\kenten{期間}%
\hfill{}{\scriptsize {\bfseries for} two days, {\bfseries for} three months, {\bfseries for} a long time \ldots}}
 \item \visible<11->{{\bfseries since} 〜は\kenten{起点}%
\hfill{}{\scriptsize {\bfseries since} yesterday, {\bfseries since} 2000, {\bfseries since}  he was a child \ldots}}
\end{itemize}
      \end{block}
\end{frame}
%%%%%%%%%%%%%%%%%%%%%%%%%%%%%%%%%%%%%%
\begin{frame}[plain]{Exercises}

{\small あたえられた日本語の意味になるよう、(~~~~~~~~)内の語句を並べ替えましょう。ただし2と4では、不足している1語を補ってください。なお、先頭の語は大文字ではじめてください}%
\hfill{\tiny 0210}\,{\scriptsize \myaudio{./audio/012_have_pp_keizoku_06.mp3}}

 \begin{enumerate}
  \item {\small 彼は先週の金曜日から病気だ。}\hfill%
( has / sick /  he / been )\,\,\,since last  Friday.\\
\visible<2->{He has been sick since last Friday.}
  \item 
{\small 彼らはこの町に住んで10年になります。}$[+1]$\\
( lived / they / this town / have / in )\,\,\,ten years.\\
\visible<3->{They have lived in this town for ten years.}
  \item 
{\small 昨日からずっと雨が降っています。}
\hfill{}( been / it / rainy / has )\,\,\,since yesterday.\\
\visible<4->{It has been rainy since yesterday.}
  \item 
{\small 私は彼女を2000年から知っている。}$[+1]$\hfill%
( her / known / I / have )\,\,\,2000.\\
\visible<5->{I have known her since 2000.}\hfill{\scriptsize 読み方two thousand \textipa{/t\'u: T\'aUznd/}}
 \end{enumerate}
\end{frame}
%%%%%%%%%%%%%%%%%%%%%%%%%%%%
\begin{frame}[plain]{継続用法の現在完了の目印}
\large

継続用法の現在完了とよくいっしょに使うfor（期間）とsince（起点）

forやsinceが継続の目印ととらえることができますね
\end{frame}
%%%%%%%%%%%%%%%%%%%%%%%%%%%%%
\section{否定文}
%%%%%%%%%%%%%%%%%%%%%%%%%%%%
\begin{frame}[plain,t]{否定文}
 \Large

%\visible<1->{I \,\myAnch{aux1}{NavyBlue}{have known} her for a long time.}
\visible<2->{I \,\myAnch{aux2}{NavyBlue}{\bfseries have} \textcolor{Maroon}{\bfseries not} \,\myAnch{aux3}{NavyBlue}{\bfseries seen} you for a long time.}

\visible<3->{I \textcolor{Maroon}{\bfseries haven't} {\bfseries seen} you for a long time.
}
\vspace{70pt}

\begin{block}<3->{Topics for Today}\small
\begin{itemize}\setbeamertemplate{items}[square]
 \item  \visible<3->{現在完了の否定は$\text{\bfseries have} + \text{\bfseries not} + \text{過去分詞}$\\
\mbox{}\hfill%
\visible<4->{{\scriptsize 主語が三人称単数のときは$\text{\bfseries has} + \text{\bfseries not} + \text{過去分詞}$}}}
 \item  \visible<5->{\textbf{haven't}\,($=\text{\bfseries have not}$)、\textbf{hasn't}\,($=\text{\bfseries has not}$) という短縮形もあります}
% \item  She\textcolor{orange}{'s not} ということもあります
\end{itemize}
      \end{block}

\vspace{-10pt}

\hfill{\tiny 0118}\,{\scriptsize \myaudio{./audio/012_have_pp_keizoku_07.mp3}}

\normalsize
\begin{textblock*}{0.4\linewidth}(360pt,50pt)
\visible<4->{\begin{tikzpicture}
\bear[
%speech={\scriptsize 現在完了},
signpost=\scalebox{0.5}{
\parbox{2.2cm}{\color{black}
\centering 会ってない\\期間が\\長いってこと}},
signcolour= brown!70!gray,
signback=white!80!brown
]
\end{tikzpicture}}
\end{textblock*}
\end{frame}
%%%%%%%%%%%%%%%%%%%%%%%%%%%
\section{疑問文と答え方}
\begin{frame}[plain,t]{疑問文と答え方}
 \Large

\mbox{}\phantom{~You~~~~~}%
\visible<1->{You \myAnch{aux1}{NavyBlue}{have} known her for a long time.}

\vspace{20pt}

\visible<2->{\myAnch{aux2}{NavyBlue}{Have} you known her for a long time\fcolorbox{NavyBlue}{white}{?}}
\hfill{\tiny 0216}\,{\scriptsize \myaudio{./audio/012_have_pp_keizoku_08.mp3}}


\mbox{}\hfill%
\visible<4->{$\left\{\begin{tabular}{l}
	 Yes, I have.\\
No, I have not($=$ haven't)
	\end{tabular}\right\}$}

\begin{tikzpicture}[remember picture,overlay]
\visible<3->{\draw[NavyBlue, line width=3pt,opacity=.5,->] (aux1.south) to[out=-90,in=90]node[sloped,above,text=black,font=\tiny,pos=.67]{先頭へ} (aux2.north);}
\end{tikzpicture}

%\vspace{5pt}

\begin{block}<5->{Topics for Today}\small
\begin{itemize}\setbeamertemplate{items}[square]
 \item \visible<6->{現在完了形の疑問文は$\text{\bfseries Have} + \text{\bfseries S} + \text{過去分詞\ldots\,?}$}\\
\mbox{}\hspace{0\zw}%
\hfill\visible<7->{{\scriptsize 主語が三人称単数のときは$\text{\bfseries Has} + \text{\bfseries S} + \text{過去分詞\ldots{\,?}}$}}
 \item \visible<8->{答え方は\textbf{Yes, S} $+$ \textbf{have}.または\textbf{No, S} $+$ \textbf{have not}($=$\textbf{haven't}).}\\
\hfill\visible<9->{{\scriptsize 主語が三人称単数のときは\textbf{Yes, S} $+$ \textbf{has}.または\textbf{No, S} $+$ \textbf{has not}($=$\textbf{hasn'}t).}}
\end{itemize}
      \end{block}

\begin{textblock*}{0.4\linewidth}(70pt,100pt)
\visible<3->{\begin{tikzpicture}
\duck[signpost=\scalebox{0.25}{
\parbox{2.5cm}{\color{black}\centering
{\LARGE Haveだけ\\前に}！}},
signcolour=brown!70!gray,
signback=white!80!brown,
graduate=gray!20!black,
tassel=red!70!black,
laughing,
%speech={\tiny 反復練習！}
]
\end{tikzpicture}}
\end{textblock*}
\end{frame}
%%%%%%%%%%%%%%%%%%%%%%%%%%%%%%%%%%%
\begin{frame}[plain]{Exercises}
 
{\small あたえられた日本語の意味になるよう、(~~~~~~~~) 内の語句を並べ替えましょう。なお、先頭の語は大文字ではじめてください}%
\hfill{\tiny 0149}\,{\scriptsize \myaudio{./audio/012_have_pp_keizoku_09.mp3}}


\begin{enumerate}
 \item {\small わたしたちは昨年からジョンにあっていません。}\\
We  ~( not / seen / have ) John since last year.\\
\visible<2->{We have not seen John since last year.}
 \item  {\small わたしは長いあいだ寿司を食べていません。\scalebox{3}{\twemoji{sushi}}}\\
I ~~( eaten /  have / not )~~sushi~~( long / a / for / time ).\\
\visible<3->{I have not eaten sushi for a long time.}
 \item  {\small あなたのおばさんは先月からパリに滞在していますか}？\scalebox{3}{\twemoji{baguette bread}}\bcdfrance\\
( stayed /  your aunt / has ) in Paris ( last / month / since ) ? \\
\visible<4->{Has your aunt stayed in Paris since last month?}
\end{enumerate}

\begin{textblock*}{0.4\linewidth}(360pt,125pt)
\visible<1->{\includegraphics[width=.4\textwidth]{./images/sushi.png}}
\end{textblock*}
\end{frame}
%%%%%%%%%%%%%%%%%%%%%%%%%%%%%%%%%%%%%%%%%
\section{How long 〜?}
%%%%%%%%%%%%%%%%%%%%%%%%%%%%%%%%%%%%%%%%%
\begin{frame}[plain]{How long 〜?}

\visible<1->{You have stayed in Paris \myAnch{long1}{Maroon}{for two weeks}.}
 
\vspace{15pt}

\hfill\visible<2->{cf. \textcolor{NavyBlue}{\bfseries Have you stayed} in Paris for two weeks?}

\vspace{15pt}

\visible<3->{\myAnch{long2}{Maroon}{\bfseries How long} \textcolor{NavyBlue}{\bfseries have you stayed} in Paris?}%
\hfill\visible<4->{{\scriptsize どれくらいの期間}}

\begin{tikzpicture}[remember picture,overlay]
 \visible<4->{\draw[Maroon,->,line width=3pt,opacity=.5] (long1.south) to[out=-90,in=90] node[sloped,above,text=black,font=\tiny,pos=.5]{How longに置き換え先頭へ} node[sloped,below,text=black,font=\tiny,pos=.5]{後ろは疑問文の語順} (long2.north);}
\end{tikzpicture}

\begin{block}<5->{Topics for Today}\small
\begin{itemize}\setbeamertemplate{items}[square]
 \item $\text{\bfseries How long have} + \text{\bfseries S} + \text{過去分詞 \ldots ?}$は「継続している期間」をたずねる表現\\
\mbox{}\hfill「どのくらいのあいだ(期間)〜ですか」\\
\hfill{\scriptsize 主語が三人称単数のときは$\text{\bfseries How long has} + \text{\bfseries S} + \text{過去分詞 \ldots ?}$}
 \item \textbf{How long}の後ろは疑問文の語順

\end{itemize}
      \end{block}

\vspace{-10pt}

\hfill{\tiny 0143}\,{\scriptsize \myaudio{./audio/012_have_pp_keizoku_10.mp3}}

\end{frame}
%%%%%%%%%%%%%%%%%%%%%%%%%%%%%%%
\begin{frame}[plain]{Exercises}

{\small あたえられた日本語の意味になるよう、(~~~~~~~~) 内の語句を並べ替えましょう。なお、先頭の語は大文字ではじめてください}%
\hfill{\tiny 0205}\,{\scriptsize \myaudio{./audio/012_have_pp_keizoku_11.mp3}}

\vspace{-5pt}

\begin{enumerate}
 \item どのくらいのあいだスペイン語を勉強してきましたか？\\
( studied / long / have / you / how ) Spanish?\\
\visible<2->{How long have you studied Spanish?}
 \item 彼女はローマにどのくらいのあいだ住んでいますか？\\
( has / lived / long / she / how ) in Rome?\\
\visible<3->{How long has she lived in Rome?}
 \item これまでどのくらいのあいだこの会社で働いてきましたか？\\
( worked /  you / long / how / have ) at this company?\\
\visible<4->{How long have you worked at this company?}
 \item ジェニファーを知ってどのくらいの期間になりますか？\\
( long / known / you / have / how ) Jennifer?\\
\visible<5->{How long have you known Jennifer?}
\end{enumerate}

\begin{textblock*}{0.4\linewidth}(386.5pt,60pt)
\visible<1->{\worldflag[width=1cm,framewidth=0.3mm,framecolor=black]{ES}}
\end{textblock*}

\begin{textblock*}{0.4\linewidth}(386.5pt,105pt)
\visible<1->{\worldflag[width=1cm,framewidth=0.3mm,framecolor=black]{IT}}
\end{textblock*}

\begin{textblock*}{0.4\linewidth}(337pt,140pt)
\visible<1->{\includegraphics[width=.6\textwidth]{./images/worker.png}}
\end{textblock*}
\end{frame}
%%%%%%%%%%%%%%%%%%%%%%
\section{聞いてみよう、読んでみよう{\tiny 0043}\,{\scriptsize \myaudio{./audio/012_have_pp_keizoku_reading.mp3}}}
%%%%%%%%%%%%%%%%%%%%%%
%%%%%%%%%%%%%%%%%%%%%%%%%%%%%%%
\begin{frame}[plain]
 
\includegraphics[width=1.01\textwidth]{./images/nanobanana-output/012_have_pp_keizoku_reading.png}

\vspace{-15pt}

\hfill{\tiny 0043}\,{\scriptsize \myaudio{./audio/012_have_pp_keizoku_reading.mp3}}

\end{frame}
%%%%%%%%%%%%%%%%%%%
\begin{frame}[plain,t]{Exercises}

\begin{tcolorbox}[colframe=ForestGreen,
  colback=ForestGreen!10!white,
  colbacktitle=ForestGreen!40!white,
  coltitle=black, %fonttitle=\bfseries,
before upper={\setlength{\parindent}{1.25em}},
 title=My Life and My Best Friend\mbox{}\hfill{\tiny 0043}\,{\scriptsize \myaudio{./audio/012_have_pp_keizoku_reading.mp3}}
]

Hello.
My name is Yumi.
I live in Chiba.
I have lived here since 2010.
I like this town very much.

I have a best friend named\footnotemark[1] Ken.
I have known him for five years.
We go to the same school.
Ken has had a dog since last month.
The dog is very cute.
But Ken has been sick since yesterday.
I haven't seen him since then.
I hope he gets well\footnotemark[2] soon.

How long have you lived in your town?

\vspace*{20pt}

\scriptsize
$^{1}$named Ken：ケンという名前の（ともだち）\hfill{}$^{2}$get well：元気になる\hfill$^{3}$soon：すぐに、まもなく

\end{tcolorbox}

\end{frame}
%%%%%%%%%%%%%%%%%%%%%%%%%%%%%%%%%
\begin{frame}[plain]{大意}
 
\begin{tcolorbox}[title=私の生活と親友]
　 こんにちは。
わたしの名前はユミです。
わたしは千葉に住んでいます。
わたしは2010年からここに住んでいます。
わたしはこの町が大好きです。

　わたしにはケンという名前の親友がいます。
わたしは彼のことを知って5年になります。
わたしたちは同じ学校にかよっています。
ケンは先月から犬を飼っています。
その犬はとてもかわいいです。
しかし、ケンはきのうから具合が悪いです。
わたしはそのときから彼に会っていません。
彼がすぐにげんきになることを願っています。

　あなたはどのくらいのあいだ、あなたの町に住んでいますか？
\end{tcolorbox}
\end{frame}
%%%%%%%%%%%%%%%%%%%%%%%%%%%%%%%%%%%%%%%%%%%%%%%%%%%%%%%%%%%%%%
\section{まとめ}
%%%%%%%%%%%%%%%%%%%%%
\begin{frame}<1-12>[plain]{まとめ}

\begin{block}<1->{「継続」を表す現在完了}
\small

\begin{itemize}\setbeamertemplate{items}[square]
 \item[]<2-> 基本: 現在完了形\,\Circled[fill color=white]{\,\textbf{have} $+$ 過去分詞\,}\,は「過去と現在にまたがる表現」

 \item \visible<3->{\<「継続」を表す用法があります\hfill{\scriptsize 過去のあるときから、いままでずっと}}
 \item \visible<4->{for（期間）やsince（起点）がセットで使われることがよくあります}
  \item  \visible<5->{否定は$\textbf{have} + \textbf{not} + \textbf{過去分詞}$%\\
%\mbox{}\hspace{9\zw}%
\hfill\visible<6->{{\scriptsize 主語が三人称単数のときは$\textbf{has} + \textbf{not} + \text{過去分詞}$}}}\\
\mbox{}\hfill\visible<7->{{\scriptsize 短縮形は\textbf{haven't}($=\textbf{have not}$)、\textbf{hasn't}($=\textbf{has not}$)}}
% \item  She\textcolor{orange}{'s not} ということもあります

 \item \visible<8->{疑問文は$\textbf{Have} + \textbf{S} + \text{過去分詞\ldots\,?}$}\hfill\visible<9->{{\scriptsize 主語が三人称単数のときは$\textbf{Has} + \textbf{S} + \text{過去分詞\ldots{\,?}}$}}
 \item \visible<10->{答え方は\textbf{Yes, S} $+$ \textbf{have}.または\textbf{No, S} $+$ \textbf{have not}($=\textbf{haven't}$).}\\
\mbox{}\hfill\visible<11->{{\scriptsize 主語が三人称単数のときは\textbf{Yes, S} $+$ \textbf{has}.または\textbf{No, S} $+$ \textbf{has not}($=\textbf{hasn't}$).}}

 \item \visible<12->{$\textbf{How long have} + \textbf{S} + \text{過去分詞 \ldots ?}$は「継続している期間」をたずねる表現}\\%
%\mbox{}\hfill\visible<7->{「どのくらいのあいだ(期間)〜ですか」　　　}\\
\mbox{}\hfill\visible<13->{{\scriptsize 主語が三人称単数のときは$\textbf{How long has} + \textbf{S} + \text{過去分詞 \ldots ?}$}}
\end{itemize}
      \end{block}
\end{frame}
%%%%%%%%%%%%%%%%%%%%%%%%%%%%%%%%
\begin{frame}[plain]

 {\tiny audio\_overview 0857}\,{\scriptsize \myaudio{./audio/overview/012_have_pp_keizoku_audio_overview.m4a}}
\end{frame}
%%%%%%%%%%%%%%%%%%
\end{document}

